\section{对联生成的结构调整与微调实验}

预训练迁移对照表明，古典语料能够改善对联监督学习的起点，但初始模型仍存在词组组织和形式对应错误。本节从这些实际问题出发介绍最终方法，随后比较原框架直接微调、增强模型无规则解码和增强模型规则约束解码，并结合最后一轮消融讨论各设计的作用。

\Needspace{11\baselineskip}
\subsection{方法与设计依据}

\textbf{任务组织与改进动机。}直接微调保留原始6层字符级GPT，将输入组织为“BOS、UP、上联、DOWN、下联、EOS”，仅监督下联与EOS。该模型已能形成下联体裁，但在相同输入上生成“且舒敞胸襟／且放放胸怀”“诗歌伴雨吟／墨韵伴风吟”，分别出现无依据重复和同位照用；多字词及其对应关系没有得到显式表示。句法增强对联研究通过词性与依存关系刻画词语间联系\cite{song}，提示应在逐字预测之外引入更明确的对应信息。本文最终采用词汇融合与位置连接，模型结构见图~\ref{fig:lexical}。

\textbf{词汇表示与注意力融合。}针对字符表示中缺少整词信息的问题，本文借鉴LEBERT的字符到词注意力，将外部词汇信息注入Transformer中间层\cite{lebert}。词库整合\href{https://github.com/fxsjy/jieba}{jieba}、\href{https://github.com/thunlp/THUOCL}{THUOCL}、\href{https://github.com/mozillazg/phrase-pinyin-data}{词语拼音表}、\href{https://github.com/mapull/chinese-dictionary}{mapull汉语词典}、\href{https://github.com/iDvel/rime-ice}{rime-ice}及\href{https://github.com/thunlp/OpenHowNet}{OpenHowNet}等资源，规范化去重后选取训练中出现的202,585条候选，每字保留至多8个重叠匹配，最长8字。保留重叠跨度使模型能够在多种候选中学习选择。

词向量采用\href{https://huggingface.co/ShannonAI/ChineseBERT-base}{ChineseBERT}的条件上下文特征。该模型融合字符、字形和拼音信息\cite{chinesebert}，为词语提供随语境变化的表示。实验冻结其参数，以降低任务适配的更新规模；对已知条件句编码后，对词跨度均值池化，经$P_c:768\to128$投影并加上可训练词条残差$r_w$，再映射至GPT的256维空间。以首层GPT字符状态$h_i$为查询，融合计算为
\begin{align}
 v_w&=\operatorname{MLP}\!\left(P_c\operatorname{Mean}(H_{s_w:e_w})+r_w\right),\label{eq:word}\\
 \alpha_{iw}&=\operatorname{softmax}_{w\in C(i)}(h_i^{\mathsf T}Av_w),\qquad
 c_i=\sum_{w\in C(i)}\alpha_{iw}v_w.\label{eq:lexical}
\end{align}
$C(i)$为覆盖位置$i$的候选集合。$c_i$经归一化、Dropout和零初始化投影后，残差注入第一层之后的条件句状态；无候选时增量为零。融合位置使后续层持续处理词汇信息。LEBERT原用于BERT序列标注，本文将其融合思想迁移到自建因果GPT；ChineseBERT作为冻结特征来源。正向只编码上联，反向只编码候选下联，均以该方向已知条件构造特征。

\begin{figure}[!htbp]
\centering\resizebox{\linewidth}{!}{% One-direction architecture. The reverse model swaps x and y, with independent weights.
\begin{tikzpicture}[x=1cm,y=1cm,
 font=\fontsize{9.6}{11.5}\selectfont,
 layer/.style={draw=black!65,line width=.5pt,rounded corners=1pt,align=center,minimum height=.44cm,inner sep=3pt},
 tok/.style={draw=black!65,fill=white,minimum width=.47cm,minimum height=.32cm,inner sep=1pt},
 ar/.style={-{Stealth[length=1.6mm]},line width=.55pt},
 plus/.style={circle,draw=black!70,fill=white,minimum size=.34cm,inner sep=0pt}]
\node[anchor=west,font=\small\bfseries] at (-.2,10.6) {(a) 条件表示与词汇增强解码器};
% Frozen condition encoder, lower left; embeddings explicitly show representation channels.
\node at (1.45,-.15) {条件句 $x_1,\ldots,x_n$};
\node[layer,fill=purple!12,minimum width=.78cm] (ce) at (.25,.60) {Char};
\node[layer,fill=orange!18,minimum width=.78cm] (ge) at (1.45,.60) {Glyph};
\node[layer,fill=modelblue!18,minimum width=.78cm] (pe) at (2.65,.60) {Pinyin};
\foreach \n in {ce,ge,pe}{\draw[ar] (1.45,.06)--(\n.south);}
\node[layer,fill=black!7,minimum width=3.2cm] (fusion) at (1.45,1.53) {Fusion $+$ Position};
\foreach \n in {ce,ge,pe}{\draw[ar] (\n.north)--(fusion.south);}
\draw[draw=graytext,fill=black!4] (-.10,2.07) rectangle (3.0,2.70);
\draw[draw=graytext,fill=black!5] (-.17,2.13) rectangle (2.93,2.76);
\node[layer,fill=black!7,minimum width=3.1cm,minimum height=.63cm] (bert) at (1.38,2.52) {ChineseBERT Encoder};
\draw[ar] (fusion)--(bert);
\draw[dashed,rounded corners=2pt,draw=graytext] (-.45,.2) rectangle (3.35,3.03);
\node[anchor=west,text=graytext,font=\fontsize{9}{10.5}\selectfont,fill=white,inner sep=1pt] at (-.27,3.03) {Frozen parameters};
% Encoder output tensor H_x.
\foreach \i in {0,...,5}{\foreach \j in {0,...,2}{\draw[fill=black!16,draw=white,line width=.3pt] (.45+.33*\i,3.40+.12*\j) rectangle ++(.32,.11);}}
\draw[ar] (bert.north)--(1.38,3.39);
\node[anchor=west] at (2.55,3.58) {$H_x$};
\node[layer,fill=orange!12,minimum width=3.45cm] (pool) at (1.45,4.35) {Span mean pooling};
\draw[ar] (1.45,3.78)--(pool);
\node[layer,fill=white,minimum width=1.0cm] (dict) at (-1.0,4.35) {$\mathcal D$};
\draw[ar] (dict)--node[above,font=\fontsize{8.8}{10}\selectfont]{spans}(pool);
\node[layer,fill=orange!14,minimum width=3.45cm] (proj) at (1.45,5.22) {$P_c\;(768\to128)$};
\node[plus] (pw) at (1.45,6.06) {$+$};
\node[layer,fill=orange!14,minimum width=1.0cm] (rw) at (-1.0,6.06) {$r_w$};
\draw[ar] (pool)--(proj);\draw[ar] (proj)--(pw);\draw[ar] (rw)--(pw);
\node[layer,fill=orange!14,minimum width=3.45cm] (mlp) at (1.45,6.90) {MLP $128\to256$};\draw[ar] (pw)--(mlp);
% Word candidate vectors.
\foreach \i in {0,...,3}{\draw[fill=orange!22,draw=black!40,line width=.3pt] (.45+.5*\i,7.57) rectangle ++(.34,.55);}
\draw[ar] (mlp)--(1.45,7.55);\node at (1.45,8.42) {$V_i=\{v_w:w\in C(i)\}$};
% Shared sequence decoder tower, right of the frozen encoder.
\foreach \x/\t in {4.35/x_1,4.98/\cdots,5.61/x_n,6.24/y_{<j}}{\node[tok] at (\x,.25) {$\t$};}
\node[anchor=west,font=\fontsize{8.8}{10}\selectfont] at (4.08,-.35) {BOS / UP / DOWN 标记省略};
\node[layer,fill=green!12,minimum width=3.18cm] (emb) at (5.30,1.18) {Token Embedding $+$ Position};
\foreach \x in {4.35,4.98,5.61,6.24}{\draw[ar] (\x,.43)--(\x,.94);}
\node[layer,fill=modelblue!18,minimum width=3.18cm,minimum height=.65cm] (b1) at (5.30,2.44) {Causal GPT Block 1};\draw[ar] (emb)--(b1);
% Source query tap and source-only additive injection.
\node[plus] (lx) at (5.30,4.15) {$+$};
\draw[ar] (b1)--node[right,inner sep=1pt] {$H^{(1)}$}(lx);
\node[layer,fill=orange!18,minimum width=3.45cm,minimum height=.64cm] (adapter) at (1.45,9.10) {Lexicon Adapter};
\draw[ar,draw=ruleorange] (1.45,8.62)--(adapter.south);
\draw[ar,draw=modelblue] (b1.east)--(7.65,2.44)--(7.65,9.75)--(1.45,9.75)--(adapter.north);
\node[anchor=west,text=modelblue,fill=white,inner sep=1pt,font=\fontsize{9}{10.5}\selectfont] at (7.65,4.9) {$Q$};
\draw[ar,draw=ruleorange] (adapter.east)--(3.35,9.10)--(3.35,4.15)--node[above,fill=white,inner sep=1pt] {$\Delta_x$}(lx.west);
\node[anchor=west,font=\fontsize{9}{10.5}\selectfont] at (5.55,4.15) {source only};
\node[layer,fill=purple!10,minimum width=3.18cm] (a1) at (5.30,5.06) {Position-aligned residual};\draw[ar] (lx)--(a1);
\draw[rounded corners=2pt,fill=modelblue!4,draw=black!55] (3.54,5.70) rectangle (7.03,7.55);
\node[layer,fill=modelblue!18,minimum width=2.98cm] (blk) at (5.30,6.19) {Causal GPT Block};
\node[layer,fill=purple!10,minimum width=2.98cm] (ali) at (5.30,7.00) {Position-aligned residual};
\draw[ar] (a1)--(blk);\draw[ar] (blk)--(ali);
\node[anchor=west,font=\fontsize{9}{10.5}\selectfont] at (7.06,6.65) {$5\times$};
\node[layer,fill=green!12,minimum width=3.18cm] (head) at (5.30,8.08) {LayerNorm $+$ Linear $E^{\mathsf T}$};\draw[ar] (ali)--(head);
\node[layer,fill=green!12,minimum width=3.18cm] (soft) at (5.30,8.95) {Softmax};\draw[ar] (head)--(soft);
\draw[ar] (soft)--(5.30,9.75);\node at (5.30,10.02) {$p_f(y_j\mid x,y_{<j})$};
% Right inset: actual attention and alignment operators, rather than prose boxes.
\node[anchor=west,font=\small\bfseries] at (9.3,10.6) {(b) 词汇适配器};
\draw[dashed,draw=black!45,rounded corners=2pt] (9.20,5.14) rectangle (14.9,10.10);
\node (q) at (10.1,5.62) {$h_i$};\node (v) at (13.4,5.62) {$V_i$};
\node[layer,fill=orange!12,minimum width=3.75cm] (score) at (12,6.48) {$e_{iw}=h_i^{\mathsf T}Av_w$};
\draw[ar] (q)--(score);\draw[ar] (v)--(score);
\node[layer,fill=orange!12,minimum width=3.75cm] (alpha) at (12,7.32) {$\alpha_{iw}=\operatorname{softmax}_w(e_{iw})$};\draw[ar] (score)--(alpha);
\node[layer,fill=orange!12,minimum width=3.75cm] (sum) at (12,8.16) {$c_i=\sum_w\alpha_{iw}v_w$};\draw[ar] (alpha)--(sum);
\draw[ar] (v.east)--(14.55,5.62)|-(sum.east);
\node[layer,fill=orange!12,minimum width=3.75cm] (delta) at (12,9.04) {LayerNorm / Dropout / $W_o$};\draw[ar] (sum)--(delta);
\draw[ar] (delta)--(12,9.70);\node[anchor=south] at (12,9.7) {$\Delta_i$};
\node[anchor=west,font=\small\bfseries] at (9.3,4.55) {(c) 同位置源端残差};
\draw[dashed,draw=black!45,rounded corners=2pt] (9.20,.28) rectangle (14.9,4.15);
\node[layer,fill=modelblue!14,minimum width=1.4cm] (sx) at (10.3,1.00) {$h_{x_j}^{(\ell)}$};
\node[layer,fill=orange!13,minimum width=1.4cm] (ty) at (13.6,1.00) {$h_{t(j)}^{(\ell)}$};
\node[layer,fill=purple!12,minimum width=1.4cm] (A) at (10.3,2.30) {$A^{(\ell)}$};\draw[ar] (sx)--(A);
\node[plus] (pa) at (13.6,2.30) {$+$};\draw[ar] (ty)--(pa);\draw[ar] (A)--(pa);
\draw[ar] (pa)--(13.6,3.35);\node[anchor=south] at (13.6,3.35) {$\widetilde h_{t(j)}^{(\ell)}$};
\node[anchor=west,text=graytext,font=\fontsize{9}{10.5}\selectfont] at (9.55,.40) {条件端\hspace{2.35cm}目标预测端};
\end{tikzpicture}
}
\caption{词汇增强对联模型。冻结ChineseBERT提供条件跨度表示，词汇适配器在GPT第一层之后注入条件端增量；每层的位置残差连接上联与目标预测位置。$\mathcal D$为词库，$Q$为字符查询，$r_w$为词条残差；蓝线传递查询，橙线传递词汇增量，虚线边界表示冻结参数。正反向主干采用独立参数。}\label{fig:lexical}
\end{figure}

\Needspace{7\baselineskip}
\textbf{位置对应与监督目标。}对联通常等长，相同位置承担对应功能。为强化这一联系，本文在每层预测下联第$j$字的状态中加入上联对应状态的投影，作为本实验的位置先验设计：
\begin{equation}
 \widetilde h_{t(j)}^{(\ell)}=h_{t(j)}^{(\ell)}+A^{(\ell)}h_{s(j)}^{(\ell)},\qquad
 \mathcal L_f=\mathcal L_{\mathrm{CE}}+0.5\mathcal L_{\mathrm{struct}}.\label{eq:align}
\end{equation}
$s(j)$为上联位置，$t(j)$为预测目标字的位置，投影以零初始化。辅助损失依据参考下联，在有效位置抑制同位照用，鼓励参考支持的重复对应，并抑制无依据重复；重复关系考察先前至多7个位置。该设计把形式关系转化为概率优化目标，与提供词语信息的词汇通道共同训练。

\textbf{数据与微调设置。}配对语料来自\href{https://github.com/wb14123/couplet-dataset}{couplet-dataset}、\href{https://github.com/v-zich/couplet-clean-dataset}{couplet-clean-dataset}和\href{https://gitlab.com/feng-7/VV-couplet}{VV-couplet}，按上联及父联关系组织划分。最终训练944,256对，验证和测试分别为51,600与49,448条配对；选点使用固定3,663对，生成评测使用2,617个锁定上联。反向学习$p_r(x\mid y)$，额外移除33,099对与保留集下联相交的训练记录。表~\ref{tab:sftconfig}汇总主要设置。

\begin{table}[H]
\centering\small
\caption{直接微调基线与最终增强模型的训练设置}\label{tab:sftconfig}
\begin{tabularx}{\linewidth}{Yccc}
\toprule 设置&原框架直接微调&增强正向模型&增强反向模型\\\midrule
训练配对数&400,290&944,256&911,157\\
采样／遍历轮数&3&3&3\\
有效批量&256&128&128\\
优化更新数&4,692&22,131&21,357\\
学习率&$3\!\times\!10^{-4}$峰值&$2\!\times\!10^{-5}$&$2\!\times\!10^{-5}$\\
训练目标&下联及EOS交叉熵&交叉熵＋形式辅助项&反向交叉熵\\\bottomrule
\end{tabularx}
\end{table}

各模型共享同源预训练起点，增强模型在已有任务权重上继续适配，以验证交叉熵选点，采用AdamW、BF16及梯度范数裁剪1。基线训练仅含五／七言；增强正向模型对其他长度赋予4倍采样权重，累计呈现2,832,768个样本。每个增强方向含34,759,428个可训练参数，冻结ChineseBERT另含147,337,984参数；反向模型独立训练并累计呈现2,733,471个样本。

\textbf{候选选择与规则补强。}为减少只依赖高频表达的候选选择，参考双向条件概率重排思想\cite{mmi}，采用含EOS的平均对数概率评分：
\begin{equation}
 S(x,y)=\frac{\log p_f(y\mid x)}{|y|+1}+0.25\frac{\log p_r(x\mid y)}{|x|+1}.\label{eq:score}
\end{equation}
反向系数由验证集选择。候选来自束宽16的束搜索及8条固定采样，采样为top-$k=20$、温度0.8。规则约束在搜索中执行汉字、等长、同位避字、重复模式及可用句脚读音检查，再按式\eqref{eq:score}选择结果。无规则组仅屏蔽协议控制词元，允许自然EOS，最多生成63个新词元；有规则组按联长结束。规则在推理阶段补强GPT的形式执行能力。

\Needspace{8\baselineskip}
\subsection{实验结果}

表~\ref{tab:maincomparison}比较相同2,617个上联及参考下联。基线采用原实验按验证损失选定的权重，补充推理使用与最终评测相同的束宽、采样数量和词元预算；其选择仅用正向分数，增强模型使用双向分数。五／七言子集含1,200条，其余长度含1,417条，分别列出以体现训练覆盖范围。

\begin{table}[htbp]
\centering\small
\caption{原框架直接微调与增强系统的生成结果}\label{tab:maincomparison}
\begin{tabularx}{\linewidth}{Yrrrr}
\toprule &\multicolumn{2}{c}{五／七言（1,200条）}&\multicolumn{2}{c}{全长度（2,617条）}\\
\cmidrule(lr){2-3}\cmidrule(lr){4-5}
系统配置&chrF&结构通过率（\%）&chrF&结构通过率（\%）\\\midrule
原框架直接微调&2.2991&68.08&1.5059&31.22\\
增强模型，无规则&3.2208&81.17&3.0937&76.88\\
增强模型，规则约束&3.2993&100.00&3.1585&100.00\\\bottomrule
\end{tabularx}
\par\vspace{3pt}\begin{minipage}{\linewidth}\footnotesize 注：训练规模、可训练参数与候选选择方式见表~\ref{tab:sftconfig}及正文。结构通过率为等长汉字、同位避字与重复模式的联合通过率。\end{minipage}
\end{table}

在五／七言子集上，原框架的等长率为100\%，结构通过率为68.08\%；增强模型无规则时分别为100\%和81.17\%。因此，常规联长中的差异主要表现为同位用字与重复关系，而全长度结果还体现了训练长度覆盖的扩展。加入规则后，两类集合均达到所定义结构检查的100\%。

\begin{table}[htbp]
\centering\small
\caption{相同上联在三个系统配置下的实际输出}\label{tab:three_stages}
\begin{tabularx}{\linewidth}{L{3.1cm}YYY}
\toprule 上联&原框架直接微调&增强模型，无规则&增强模型，规则约束\\\midrule
且舒敞胸襟&且放放胸怀&且放开眼界&更放开眼界\\
经春秋冬夏&入夏夏夏秋&历春夏秋冬&历南北东西\\
菊梦依稀秋意远&梅心寂寞月光寒&枫情依旧月光寒&梅魂寂寞月光寒\\
\bottomrule
\end{tabularx}
\end{table}


\Needspace{7\baselineskip}
表~\ref{tab:three_stages}展示具体变化。“且舒敞胸襟”的基线下联出现“放放”；增强模型改为“且放开眼界”，规则进一步避开首位的“且”。“经春秋冬夏”由连续重复“夏”转为完整枚举，再通过规则改成与季节并列结构相应的方位枚举。“菊梦依稀秋意远”的基线已较通顺，而增强无规则结果照用了第三位“依”，呈现了个体样例中的不同变化。

为分析各通道，表~\ref{tab:ablation}保留最后一轮7种条件、各两种选择的全部14组结果。各组固定最终权重；同一行正向与双向选择共用候选，屏蔽仅作用于正向模型，反向模型始终完整。除第一行外均无规则约束。

\begin{table}[htbp]
\centering\small
\caption{最后一轮固定权重推理消融（2,617个上联）}\label{tab:ablation}
\begin{tabularx}{\linewidth}{Yrrrr}
\toprule &\multicolumn{2}{c}{chrF}&\multicolumn{2}{c}{结构通过率（\%）}\\
\cmidrule(lr){2-3}\cmidrule(lr){4-5}
推理条件&正向选择&双向选择&正向选择&双向选择\\\midrule
完整模型，规则约束&3.0424&3.1585&100.00&100.00\\
完整模型，无规则&2.9625&3.0937&76.73&76.88\\
屏蔽词汇融合&2.5367&2.7158&72.18&73.02\\
屏蔽上下文投影&2.5989&2.8462&77.26&77.88\\
屏蔽词条残差&2.7005&2.9932&76.88&77.03\\
仅保留8192词匹配&2.8783&3.0386&75.54&75.12\\
屏蔽同位置源端残差&0.5263&0.6045&0.15&0.73\\\bottomrule
\end{tabularx}
\end{table}

\FloatBarrier
\subsection{结果讨论}

词汇融合为模型提供了可利用的条件信息。屏蔽该通道后，双向chrF由3.0937降至2.7158；按上联进行1,000次配对重采样，屏蔽相对完整模型的逐句chrF均值差95\%区间为$[-0.5314,-0.0638]$。这一结果与词汇信息参与上下联匹配的设计目标一致。同位置残差屏蔽后约24\%输出为空，显示当前权重的条件生成依赖该直接通道。

规则约束的主要贡献是形式执行。双向结构通过率由76.88\%升至100\%，相差23.12个百分点；chrF由3.0937变为3.1585。前者来自搜索可行域中的显式检查，后者反映候选文本变化。对照文本显示，规则能够排除特定复制和重复错误，词汇及位置通道则参与候选内容的形成，两者作用于不同环节。

偏旁技巧、典故理解与整联内容关系见第\ref{sec:limitations}节的案例分析。
